% =============================================================================
% Chapter 1: The Causal Boundary — Physical AI Lecture Slides
% =============================================================================
% Built around the chapter thesis:
%   "Once a command becomes physical work it cannot be recalled,
%    so delay costs distance."
% One running case (the warehouse mobile manipulator at its 1.5 m/s drive
% limit, a worker stepping out at a rack end 1.10 m ahead) carries the deck
% from the causal boundary through the stopping budget (133.6 ms of delay,
% 762.9 mm of stop) to the proposal boundary, the Tempe collision, and the
% four bedrock laws that organize the volume.
%
% Slides are written to be read without a speaker: each frame states its
% claim in the title, shows the evidence, and closes with a bottom line.
% Speaker notes add pacing, questions, and misconceptions.
%
% Numbers are the rendered chapter's computed values (LEGO cells backed by
% MLSysIM) at the commit this deck was built against.
% =============================================================================
\documentclass[aspectratio=169, 12pt]{beamer}
\usepackage{../../assets/beamerthememlsys}

\mlsyssetup{
  volume       = {Volume IV},
  chapter      = {Chapter 1},
  logo         = {../../assets/img/logo-mlsysbook.png},
  instlogo     = {../../assets/img/logo-harvard.png},
  chaptertitle = {The Causal Boundary},
}

% --- Fonts ---
\usepackage[T1]{fontenc}
\usepackage[scaled=0.9]{helvet}
\usepackage{courier}
\renewcommand{\familydefault}{\sfdefault}

% --- Packages ---
\usepackage{booktabs}
\usepackage{amsmath}
\usepackage{array}
\newcolumntype{L}[1]{>{\raggedright\arraybackslash}p{#1}}

% --- Image paths ---
\graphicspath{
  {images/}
}

% --- Chapter-specific macros ---
\newcommand{\taudelay}{\tau_{\text{delay}}}
\newcommand{\abrake}{a_{\text{brake}}}
\newcommand{\dstop}{d_{\text{stop}}}
\newcommand{\Dclear}{D_{\text{clear}}}
% Bottom line: the one sentence a reader without a speaker should keep.
\newcommand{\bottomline}[1]{%
  \vfill
  \begin{mlsyscard}[top=0.15em,bottom=0.15em,fontupper=\footnotesize]{crimson}\textbf{Bottom line.} #1\end{mlsyscard}%
}
\definecolor{bookblue}{HTML}{006395}
\definecolor{bookorange}{HTML}{E67817}
\definecolor{bookgreen}{HTML}{008F45}
\definecolor{bookviolet}{HTML}{7E317B}
\definecolor{bookred}{HTML}{CB202D}
\definecolor{bookbrown}{HTML}{78492A}

% Compact reference entry for a long reference list
\newcommand{\mlsysrefsmall}[2]{{\footnotesize\textcolor{crimson}{\textbf{[#1]}} #2}\\[0.01cm]}

% --- Helper: safe image include ---
\newcommand{\safeimg}[2][width=\textwidth,keepaspectratio]{%
  \IfFileExists{images/#2}{\includegraphics[#1]{#2}}{%
    \IfFileExists{#2}{\includegraphics[#1]{#2}}{%
      \fbox{\parbox[c][2.5cm][c]{0.85\linewidth}{%
        \centering\footnotesize\textcolor{midgray}{[Missing image]}}}%
    }%
  }%
}

% --- Section count (must match actual \section{} count) ---
\setcounter{mlsystotalsections}{8}

\title{The Causal Boundary}
\author{Vijay Janapa Reddi}
\institute{Harvard University}
\date{}

\begin{document}

% =============================================================================
% TITLE SLIDE
% =============================================================================
\mlsystitle{The Causal Boundary}{Once a command becomes physical work it cannot be recalled, so delay costs distance.}{cover_boundary_blueprint.png}

% =============================================================================
% LEARNING OBJECTIVES
% =============================================================================
\begin{frame}{Learning Objectives}
\note{
[1 min]
% -- NARRATE
Read the subtitle aloud: it is the lecture in one line. The first half
(cannot be recalled) is irreversibility; the second half (delay costs
distance) is the stopping budget. Every objective below is one of those
two ideas applied to a different question.
% -- FLEX
[CORE] Do not dwell; return to this list at the wrap-up.
}

\small
By the end of this lecture you should be able to:
\begin{enumerate}\setlength{\itemsep}{0pt}
  \item \textbf{Trace} how learning moved from prediction to agency to embodiment, and why that gives data, compute, evaluation, and deployment physical stakes.
  \item \textbf{Explain} why an action across the \textbf{causal boundary} has consequences rollback cannot undo.
  \item \textbf{Calculate} how response delay consumes clearance with the \textbf{elementary stopping budget}.
  \item \textbf{Apply} the \textbf{scope test} to separate physical AI from advisory learning and classical control.
  \item \textbf{Distinguish} learned proposals from permission in the \textbf{five-level machine}.
  \item \textbf{Diagnose} how perception, timing, and missing redundancy interacted at Tempe.
  \item \textbf{Relate} the \textbf{four bedrock laws} to their established names and the book's four parts.
\end{enumerate}
\end{frame}

% =============================================================================
% SECTION 1: BEYOND GLASS
% =============================================================================
\section{Beyond Glass}

\begin{frame}{A better prediction can still produce a worse action}
\note{
[3 min]
% -- NARRATE
Set the scene that recurs all term. The loaded base runs down an aisle at
its 1.5 m/s drive limit. A worker steps out at the rack end 1.10 m ahead.
The model can be right about the worker and still be late: the base keeps
moving while the model computes. Every millisecond is 1.5 mm.
% -- ENGAGE
Ask: ``The model takes 100 ms longer but is 2\% more accurate. Is that a
better model for this machine?'' [Expected: you cannot answer without the
stopping budget; 100 ms is 150 mm of travel.]
% -- WARN
Students from ML assume accuracy is the quantity to maximize. Here a late
correct answer can be worse than an early adequate one.
% -- FLEX
[CORE]
}

\centering
\safeimg[width=0.95\textwidth,height=3.2cm,keepaspectratio]{aisle-scene.pdf}

\vspace{0.1cm}
\raggedright\small
The model sees the worker and proposes a stop. While it computes, the base keeps moving: at 1.5 m/s, \textbf{each millisecond of delay is 1.5 mm of travel}. The computer can still discard a proposal; it cannot take back the distance already covered.

\bottomline{Choosing an action requires reasoning about the prediction and about the physical state in which it will execute.}
\end{frame}

\begin{frame}{Three regimes differ in time, coupling, and recovery}
\note{
[3 min]
% -- NARRATE
Read the table row by row, then column by column. The advisory case is the
ML most students know: a static dataset, batched serving, and a person or
program downstream who can ignore a bad output. Agentic ML closes the loop,
but in software. Physical AI closes it through physics.
% -- ENGAGE
Ask: ``Which column changes between agentic ML and physical AI but not
between advisory and agentic?'' [Failure and recovery: rollback exists in
both digital cases.]
% -- WARN
The advisory row is not ``harmless'': a person can act on bad guidance.
The point is that the prediction itself does not command physical work.
% -- FLEX
[CORE]
}

\footnotesize
\begin{tabular}{@{}L{0.17\textwidth}L{0.22\textwidth}L{0.24\textwidth}L{0.28\textwidth}@{}}
  \toprule
  \textbf{Regime} & \textbf{Coupling} & \textbf{Time} & \textbf{Failure and recovery} \\
  \midrule
  Advisory prediction & open loop; data exogenous & batchable serving & discard or recompute before anyone acts \\
  Agentic ML & closed loop in software & elastic virtual time in a controlled environment & exception, container restart, database rollback \\
  \textbf{Physical AI} & \textbf{closed loop in physics} & \textbf{continuous; delay is distance} & \textbf{irreversible work, dissipated energy, damage} \\
  \bottomrule
\end{tabular}

\vspace{0.2cm}
\small
Output moves from a label or token, to an API payload, to \textbf{current and torque}. Authority moves from ``a person decides,'' to sandboxed scopes, to \textbf{a separation of proposal from permission}.

\bottomline{Physical AI diverges from the other two regimes in how computation, time, and consequence relate, not in the model it runs.}
\end{frame}

\begin{frame}{Embodiment removes the privileges digital ML relies on}
\note{
[3 min]
% -- LINK
Students who built agents know these privileges well, even if they never
named them.
% -- NARRATE
Left column: what a controlled software environment grants. Right column:
what the plant imposes instead. Agentic systems also get a fourth
privilege, discrete symbolic interfaces whose syntax a compiler or linter
can check before execution; a torque has no such grammar.
% -- ENGAGE
Ask: ``Which right-hand row makes the stopping budget necessary?'' [The
first: the plant does not wait, so delay becomes distance.]
% -- WARN
``Batch of one'' is not a performance tuning choice. Each proposal depends
on observations shaped by the previous motion, so there is nothing to
batch.
% -- FLEX
[CORE]
}

\small
\begin{tabular}{@{}L{0.42\textwidth}L{0.50\textwidth}@{}}
  \toprule
  \textbf{Behind the glass} & \textbf{At the causal boundary} \\
  \midrule
  \textbf{Elastic virtual time:} the environment waits while the model thinks, 200 ms or 30 s. & \textbf{Continuous time:} the plant does not pause; a machine at speed $v$ covers $v\,\taudelay$ before braking. \\
  \textbf{Transactional rollback:} exceptions, database rollback, container restart. & \textbf{Irreversible work:} Joule heat $I^2R$ and kinetic energy $\tfrac{1}{2}mv^2$ cannot be garbage-collected. \\
  \textbf{Batchable execution:} serving amortizes memory traffic over $B \gg 1$. & \textbf{Batch of one, hard deadlines:} current loops at tens of kHz, checks near 1 kHz, policies at tens of Hz. \\
  \bottomrule
\end{tabular}

\bottomline{A missed deadline on a robot does not degrade user experience; it can destabilize the plant and deliver energy that will not roll back.}
\end{frame}

\begin{frame}{Step by step, the learned output moved to the actuator}
\note{
[3 min]
% -- NARRATE
Read the rows top to bottom as eighty years. Unimate replayed commands
from a magnetic drum; Shakey planned, slowly, while the world waited.
ALVINN and Stanley learned what the world was like, and hand-written
planners decided what to do. Then networks mapped camera images straight
to motor torques, and imitation policies predicted action chunks. Now a
pretrained vision-language backbone proposes the motion, split across
rates: a slow model reads the scene, a fast one emits chunks.
% -- ENGAGE
Ask: ``At each row, what guarantee used to come from separation or a
person, and now must live inside the machine?'' [Keeping people out of
the arm's reach; a human reviewing the plan.]
% -- WARN
Do not read the bottom row as ``end-to-end learned control replaces
controllers.'' Controllers still execute the chunks beneath the proposer.
% -- FLEX
[CORE] The two research lanes (robotics learning to act, ML learning to
predict) are in the chapter's timeline figure if students ask.
}

\centering
\safeimg[width=0.95\textwidth,height=4.0cm,keepaspectratio]{learned-output-eras.pdf}

\raggedright\footnotesize
Blue: learned; gray: hand-written. The running machine follows the last row: a slow \textbf{intent model} chooses the goal (2 Hz), and a fast \textbf{chunk policy} emits short sequences of setpoints toward it.

\bottomline{Once a general-purpose model proposes motion, deciding what may permit that motion becomes a design question.}
\end{frame}

\begin{frame}{Robot policies train on a sliver of frontier compute}
\note{
[2 min]
% -- NARRATE
Gray: language models in the Epoch AI database. Green: the 25 robotics
models that carry a compute estimate at all. OpenVLA (2024) used about
350 times less training compute than the largest language model of its
year; GR00T N1 (2025) about 7,000 times less. The open squares are video
world models, which generate video rather than actions.
% -- ENGAGE
Ask: ``Why not just scale robot policies like language models?'' [Data
must be collected from machines acting in the world, and the model must
run on the machine itself.]
% -- WARN
Values are Epoch's estimates, and coverage is thin. The claim is orders
of magnitude, not exact ratios.
% -- FLEX
[OPTIONAL]
}

\begin{columns}[T]
\begin{column}{0.56\textwidth}
\safeimg[width=\textwidth,height=5.0cm,keepaspectratio]{compute-by-domain.pdf}
\end{column}
\begin{column}{0.41\textwidth}
\small
Training compute recorded by the Epoch AI database.

\vspace{0.15cm}
\textbf{OpenVLA (2024):} about \textbf{350$\times$} less than the largest language model of its year.

\vspace{0.1cm}
\textbf{GR00T N1 (2025):} about \textbf{7,000$\times$} less.

\vspace{0.1cm}
Only 25 robotics models in the database carry a compute estimate.
\end{column}
\end{columns}

\bottomline{Physical AI is not a smaller copy of the language-model scaling story: its proposers are bounded by data from machines acting and by compute the machine can carry.}
\end{frame}

\begin{frame}{Every ML systems concern acquires physical stakes}
\note{
[2 min]
% -- NARRATE
This is the first learning objective in one table. Each row is a concern
students know from ML systems, and each changes character once the output
is an irreversible action. These four rows foreshadow the book: data and
evaluation are Part II, compute and timing Part I and III, deployment
Part IV.
% -- ENGAGE
Ask: ``Which row would a standard offline benchmark miss?'' [Evaluation:
competence shows only in the closed loop.]
% -- FLEX
[CORE]
}

\small
\begin{tabular}{@{}L{0.16\textwidth}L{0.33\textwidth}L{0.43\textwidth}@{}}
  \toprule
  \textbf{Concern} & \textbf{Digital ML} & \textbf{Once outputs move mass} \\
  \midrule
  Data & a fixed corpus & the machine's own actions choose what it observes next \\
  Compute & throughput and cost & a late answer is paid for in stopping distance \\
  Evaluation & a score on held-out examples & competence that matters shows only in the closed loop \\
  Deployment & a version rollout & a release grants authority over motion near people \\
  \bottomrule
\end{tabular}

\bottomline{When a learned output becomes an irreversible action, none of these concerns can be settled inside the model; they meet where commands become physical work.}
\end{frame}

% =============================================================================
% SECTION 2: CAUSAL BOUNDARY
% =============================================================================
\section{Causal Boundary}

\begin{frame}{A tensor is a proposal; current in a winding is work}
\note{
[3 min]
% -- NARRATE
Trace a command left to right. The policy emits a floating-point tensor:
an uncommitted proposal. The permission path checks it and admits it. A
position loop compares each setpoint with the measured joint; a current
controller sets the current the correction needs. Software writes a PWM
compare register, which sets how long each power switch stays on; the
inverter turns the DC supply into current in the stator windings, and the
field produces torque on the rotor, gears, and payload.
% -- ENGAGE
Ask: ``Where exactly does the cost of changing your mind jump?'' [At the
dashed line: before it, rewriting a setpoint in memory is nearly free;
after it, a change needs physical work against inertia.]
% -- WARN
Students place the boundary at the model output. The model output is the
cheapest point in the chain to change.
% -- FLEX
[CORE]
}

\centering
\safeimg[width=0.97\textwidth,height=3.6cm,keepaspectratio]{command-chain.pdf}

\vspace{0.1cm}
\raggedright\small
A subsequent command can brake or reverse the drive current, but that correction is itself motion: it needs its own distance and must obey actuator limits.

\bottomline{The causal boundary sits where the command chain first energizes a winding; upstream of it a proposal can be discarded for free.}
\end{frame}

\begin{frame}{Definition: the causal boundary}
\note{
[2 min]
% -- NARRATE
Read the definition, then the three cards. The pitfall card is the one
students will meet in practice: killing the policy process does not stop
the base. Something must still command a stop.
% -- ENGAGE
Ask: ``A ROS node crashes mid-motion. What is the base doing?'' [Whatever
the last admitted command told it, until something else intervenes.]
% -- FLEX
[CORE]
}

\begin{mlsyscard}{crimson}
\normalsize The \textbf{causal boundary} is the interface through which a digital command produces physical forces, currents, or flows.
\end{mlsyscard}

\vspace{0.02cm}
\small
\begin{columns}[T]
\begin{column}{0.31\textwidth}
\begin{mlsyscard}{bookgreen}
\footnotesize\raggedright\textbf{Significance.} A proposal can be rejected before it changes the plant; once executed, it dissipates energy and changes later observations.
\end{mlsyscard}
\end{column}
\begin{column}{0.31\textwidth}
\begin{mlsyscard}{bookblue}
\footnotesize\raggedright\textbf{Distinction.} Restoring a command in memory does not restore the physical state. Recovery needs another physical action.
\end{mlsyscard}
\end{column}
\begin{column}{0.31\textwidth}
\begin{mlsyscard}{bookred}
\footnotesize\raggedright\textbf{Pitfall.} Assuming that canceling a software task stops its motion. The machine needs a feasible stop even when the proposer halts.
\end{mlsyscard}
\end{column}
\end{columns}

\bottomline{Irreversibility is what the causal boundary imposes: work delivered across it cannot be recalled.}
\end{frame}

\begin{frame}{An acting machine chooses its own next observation}
\note{
[3 min]
% -- NARRATE
Top: advisory output stops at an interface where a person or filter
decides. Bottom: the loop closes through the permission path, the
actuators, and the plant. The next state is $f(s_t, a_t)$, and the next
observation is drawn from that state, so the model's own actions shape
its next inputs.
% -- ENGAGE
Ask: ``ALVINN's training drivers always stayed centered. What happened
when the van drifted?'' [It veered off: the logs held no recovery
examples. Pomerleau synthesized shifted images with corrective steering.]
% -- WARN
The quadratic growth of error with horizon is a preview; Physical Data
owns the mechanics. Give the one-line idea only.
% -- FLEX
[CORE]
}

\begin{columns}[T]
\begin{column}{0.54\textwidth}
\safeimg[width=\textwidth,height=4.8cm,keepaspectratio]{causal-loop.pdf}
\end{column}
\begin{column}{0.43\textwidth}
\small
\[
s_{t+1} = f(s_t, a_t), \qquad o_{t+1} \sim p(o \mid s_{t+1})
\]
A flawed action can steer the machine into states its training data never covered, where its outputs degrade further.

\vspace{0.1cm}
Imitation learning calls this \textbf{covariate shift}: the cost of mistakes can grow quadratically with the task horizon.
\end{column}
\end{columns}

\bottomline{Offline accuracy cannot certify a closed-loop machine; long-horizon behavior has to be evaluated in the loop.}
\end{frame}

\begin{frame}{Recovery is a new constrained motion, not a rollback}
\note{
[3 min]
% -- NARRATE
A database abort restores the state it started from. A gripper that has
permanently deformed a fixture, or a winding whose insulation has
degraded from overheating, has no earlier state to return to. Every
correction is motion, bounded by torque, friction, and drive limits, and
which limit binds depends on what the machine is doing.
% -- ENGAGE
Ask: ``Why must the check happen before admission rather than after?''
[After admission, the state that would need recovery may already have no
feasible recovery.]
% -- WARN
``Reverse the motor'' can return an arm to its starting coordinates; it
cannot mend cracked material or recover energy lost to friction and heat.
% -- FLEX
[CORE]
}

\begin{columns}[T]
\begin{column}{0.47\textwidth}
\begin{mlsyscard}{bookblue}
\small\textbf{Software recovery}\\[2pt]
\footnotesize An aborted transaction restores its starting state. A restarted container begins from a clean snapshot. Recovery costs time, not physics.
\end{mlsyscard}
\end{column}
\begin{column}{0.49\textwidth}
\begin{mlsyscard}{bookbrown}
\small\textbf{Physical recovery}\\[2pt]
\footnotesize Braking, releasing a grasp, or diverting is itself motion, limited by actuator torque, friction, and drive speed and acceleration.
\end{mlsyscard}
\end{column}
\end{columns}

\vspace{0.2cm}
\small
\textbf{Which limit governs recovery:} for the moving base, stopping clearance under momentum; for the arm gripping a fragile part or bracing on a fixture, contact force, transmission compliance, and winding heat.

\bottomline{Before it admits an action, the permission path must confirm that a feasible recovery from the resulting state will still exist.}
\end{frame}

\begin{frame}{One running machine carries every chapter}
\note{
[3 min]
% -- NARRATE
Introduce the machine card. Every later chapter reads these values rather
than restating them. Point at the category column: illustrative and
chosen values are design inputs, not measurements of a real chassis.
Today we use five rows: drive limit, credible deceleration, clear
distance, pre-brake delay, and the chunk lease inside it.
% -- ENGAGE
Ask: ``Why is the aisle speed 1.3 m/s when the drive limit is 1.5?''
[Hold that question; the stopping budget answers it.]
% -- WARN
The 2 m/s$^2$ deceleration rests on a floor premise: a friction
coefficient of at least 0.204 (or 0.250 at the share the drive motors
allow). It is not a property of the brakes alone.
% -- FLEX
[CORE]
}

\begin{columns}[T]
\begin{column}{0.55\textwidth}
\footnotesize\renewcommand{\arraystretch}{0.86}
\begin{tabular}{@{}lrl@{}}
  \toprule
  \textbf{Parameter} & \textbf{Value} & \textbf{Category} \\
  \midrule
  Loaded mass & 368 kg & derived \\
  \textbf{Drive limit} & \textbf{1.5 m/s} & illustrative \\
  Aisle speed & 1.3 m/s & chosen \\
  \textbf{Credible deceleration} & \textbf{2 m/s$^2$} & illustrative \\
  \textbf{Rack-end clear distance} & \textbf{1.10 m} & illustrative \\
  \textbf{Pre-brake delay} & \textbf{133.6 ms} & derived \\
  \textbf{Chunk lease} & \textbf{60 ms} & chosen \\
  Permission tick & 1 ms & derived \\
  Intent-model rate & 2 Hz & chosen \\
  In-loop takeover & 605 ms & derived \\
  \bottomrule
\end{tabular}
\end{column}
\begin{column}{0.41\textwidth}
\small
A \textbf{warehouse mobile manipulator}: wheeled base, articulated arm, navigation lidar, wrist cameras.

\vspace{0.15cm}
It collects objects from a conveyor, opens a spring-latched cage door, and carries loads through aisles shared with people.

\vspace{0.15cm}
Ten of the card's fifteen rows; bold rows enter today's stopping budget.
\end{column}
\end{columns}

\bottomline{The machine is an illustrative engineering design; every value carries its evidence category so later chapters know what is measured and what is assumed.}
\end{frame}

% =============================================================================
% SECTION 3: STOPPING BUDGET
% =============================================================================
\section{Stopping Budget}

\mlsysfocus{The Elementary Stopping Budget}{%
$\dstop = \underbrace{v\,\taudelay}_{\text{travel during delay}} + \underbrace{\dfrac{v^2}{2\,\abrake}}_{\text{braking travel}} \;\le\; \Dclear$\\[0.4cm]
\normalsize speed $v$ \;$\cdot$\; pre-brake delay $\taudelay$ \;$\cdot$\; credible deceleration $\abrake$ \;$\cdot$\; clear distance $\Dclear$%
}

\begin{frame}{The pre-brake delay has five terms; the lease is largest}
\note{
[4 min]
% -- NARRATE
Walk the bar left to right. Observation age is 51.6 ms from the camera
exposure that first shows the worker to the handoff, and the learned
perception backbone is only 22 ms of that. Then the chunk lease, 60 ms:
the longest an admitted chunk may drive before the permission path must
act. One permission tick, one bus cycle, and 20 ms of brake onset.
% -- ENGAGE
Ask: ``The permission path re-checks every millisecond. Why does a 60 ms
lease belong in the worst case?'' [Its check sees only freshness, lease,
and kinematic bounds against the encoders. It has no camera. If the
proposer stalls, the last admitted chunk drives until the lease expires.]
% -- WARN
The budget excludes the dispatch-to-admission interval; Workload
Placement allocates it from the spare.
% -- FLEX
[CORE]
}

\centering
\safeimg[width=0.97\textwidth,height=3.6cm,keepaspectratio]{delay-waterfall.pdf}

\raggedright\footnotesize
$\taudelay = 51.6 + 60 + 1 + 1 + 20 = 133.6$ ms. The permission path has no camera: news of the worker reaches the brakes only through learned perception and a new proposal. If the proposer stalls, the last admitted chunk drives until its lease expires.

\bottomline{The learned backbone is 22 ms of the 133.6 ms delay; the lease costs more.}
\end{frame}

\begin{frame}{Your turn: does the base stop inside 1.10 m?}
\note{
[4 min] Exercise.
% -- NARRATE
Give two minutes. Students should compute the delay term, the braking
term, the sum, and the spare. Then the two follow-ups: the 22 ms backbone
is 33 mm, the 60 ms lease 90 mm.
% -- ENGAGE
Follow-up: ``By what fraction does the delay lengthen the stop beyond
braking alone?'' [200.4 / 562.5, about 35.6\%.]
% -- WARN
Students forget to convert milliseconds to seconds and get 200 m.
% -- FLEX
[CORE]
}

\small
\textbf{Given (machine card):} $v = 1.5$ m/s, $\abrake = 2$ m/s$^2$, $\Dclear = 1.10$ m, $\taudelay = 133.6$ ms.
\begin{enumerate}\setlength{\itemsep}{1pt}
  \item Does the elementary stop fit, and with how much to spare?
  \item Which term dominates?
  \item How much travel do the 22 ms backbone and the 60 ms lease each cost?
\end{enumerate}

\pause
\begin{mlsyscard}{crimson}
\small
$v\,\taudelay = 1.5 \times 0.1336 = 200.4$ mm \qquad $v^2/(2\abrake) = 1.5^2/4 = 562.5$ mm\\[3pt]
$\dstop = 762.9$ mm $\le 1{,}100$ mm: fits with \textbf{337.1 mm} to spare. Braking is 2.8$\times$ the delay travel.\\[3pt]
Backbone: $1.5 \times 22 = $ \textbf{33 mm}. Lease: $1.5 \times 60 = $ \textbf{90 mm}. Every millisecond anywhere in the chain costs 1.5 mm.
\end{mlsyscard}
\end{frame}

\begin{frame}{Delay grows with $v$, braking with $v^2$; they cross at 0.53 m/s}
\note{
[3 min]
% -- NARRATE
Equate the two terms: $v\,\taudelay = v^2/(2\abrake)$ gives the crossover
speed $v^* = 2\abrake\taudelay = 2 \times 2 \times 0.1336 \approx 0.53$
m/s. Below it, latency is the larger term; above it, braking is. The
running machine's drive limit is almost three times $v^*$.
% -- ENGAGE
Ask: ``At the drive limit, which lever saves more: 20 ms less delay, or
0.1 m/s less speed?'' [20 ms saves 30 mm; dropping to 1.4 m/s saves about
13 mm of delay and 72 mm of braking. Work it out rather than trusting
term dominance.]
% -- WARN
Dominance tells you which term is bigger, not which change saves more;
that depends on the reductions actually available.
% -- FLEX
[CORE]
}

\begin{columns}[T]
\begin{column}{0.55\textwidth}
\safeimg[width=\textwidth,height=5.0cm,keepaspectratio]{stopping-vs-speed.pdf}
\end{column}
\begin{column}{0.42\textwidth}
\small
\textbf{Crossover speed}
\[
v^* = 2\,\abrake\,\taudelay \approx 0.53 \text{ m/s}
\]
\textbf{Two levers:} shortening sensing, computation, and signaling shrinks the linear term; lowering speed shrinks both, braking quadratically.

\vspace{0.1cm}
At 1.5 m/s the delay adds \textbf{35.6\%} to the braking distance.
\end{column}
\end{columns}

\bottomline{Selecting a model on offline accuracy alone ignores the first lever: added inference latency is added travel.}
\end{frame}

\begin{frame}{Predict: what does halving the speed buy?}
\note{
[3 min] Predict, then reveal.
% -- NARRATE
Give 60 seconds. Most students halve everything. The delay term is linear,
so it halves; the braking term is quadratic, so it quarters.
% -- ENGAGE
Follow-up: ``At 0.75 m/s, which term dominates?'' [Braking, barely: 0.75
is still above $v^* \approx 0.53$.]
% -- FLEX
[CORE]
}

\centering
\vspace{0.3cm}
{\large At 1.5 m/s the stop is \textbf{200.4 mm} of delay travel + \textbf{562.5 mm} of braking = \textbf{762.9 mm}.}\\[0.35cm]
{\normalsize The base slows to \textbf{0.75 m/s}. How long is the stop now?}\\[0.2cm]
{\small\textcolor{midgray}{Write a number and say which term shrank more. 60 seconds.}}

\pause
\vspace{0.4cm}
\raggedright
\begin{mlsyscard}{crimson}
\small Delay travel halves to \textbf{100.2 mm}; braking quarters to \textbf{140.6 mm}. The stop falls from 762.9 mm to \textbf{240.8 mm}, less than a third, because speed enters both terms and the braking term quadratically.
\end{mlsyscard}
\end{frame}

\begin{frame}{Your turn: the chunk lease doubles to 120 ms}
\note{
[3 min] Exercise (the chapter's checkpoint question).
% -- NARRATE
The lease is the largest delay term, and a designer is tempted to lengthen
it so the proposer can renew less often. Each extra millisecond is 1.5 mm.
The arithmetic: 60 more ms is 90 mm, so 762.9 + 90 = 852.9 mm, which still
fits with 247.1 mm to spare in the elementary budget.
% -- WARN
``Still fits'' is only the elementary budget. The spare is not yet a
safety margin, as the next frame shows.
% -- FLEX
[OPTIONAL] Skip if the previous exercise ran long.
}

\centering
\vspace{0.3cm}
{\large Same base, same 1.5 m/s. The chunk lease grows from \textbf{60 ms} to \textbf{120 ms}.}\\[0.35cm]
{\normalsize How far does the elementary stop grow, and does it still fit inside 1.10 m?}\\[0.2cm]
{\small\textcolor{midgray}{One line of arithmetic. 60 seconds.}}

\pause
\vspace{0.4cm}
\raggedright
\begin{mlsyscard}{crimson}
\small $\taudelay = 133.6 + 60 = 193.6$ ms. The extra 60 ms costs $1.5 \times 60 = 90$ mm, so $\dstop = 762.9 + 90 = 852.9$ mm. It still fits the elementary budget, with 247.1 mm left for every allowance not yet counted.
\end{mlsyscard}
\end{frame}

\begin{frame}{The decisive delay belongs to the path that triggers the stop}
\note{
[3 min]
% -- NARRATE
Two points. First, the 337.1 mm spare is an arithmetic difference between
provisional figures. Localization error, a protective margin, tracking
error, and the longer stop of a jerk-limited brake together exceed the
clear distance at the drive limit; that is why the card sets the aisle
speed at 1.3 m/s. Second, a separate protective path with its own
sensing, timing, power, and fallback, such as a certified safety scanner
wired to the drives, bypasses perception dispatch and the lease entirely.
% -- ENGAGE
Ask: ``With a certified scanner in place, which terms of today's delay
disappear?'' [The backbone and the lease; the scanner's own sensing and
switching latency replace them. Braking travel and the allowances stay.]
% -- WARN
A pedestrian walking toward the robot shortens the gap during both phases;
standards add an approach-speed allowance.
% -- FLEX
[CORE]
}

\begin{columns}[T]
\begin{column}{0.48\textwidth}
\begin{mlsyscard}{bookred}
\small\textbf{The spare is not a safety margin}\\[2pt]
\footnotesize Later chapters add localization error, a protective margin, tracking error, and the longer jerk-limited stop. Together they exceed 1.10 m at the 1.5 m/s drive limit, which is why the card sets the \textbf{aisle speed at 1.3 m/s}. A walking pedestrian adds an approach-speed allowance.
\end{mlsyscard}
\end{column}
\begin{column}{0.48\textwidth}
\begin{mlsyscard}{bookgreen}
\small\textbf{An independent protective path}\\[2pt]
\footnotesize A certified safety scanner wired to the drives bypasses perception and the lease. Its own sensing and switching latency replace them; braking travel stays. Industrial AGVs do this today under ISO 3691-4.
\end{mlsyscard}
\end{column}
\end{columns}

\bottomline{Under a decoupled architecture, the protective path's latency, not the neural network's execution time, governs the protective stop.}
\end{frame}

\begin{frame}{Stopping clearance is one of three margins a machine spends}
\note{
[2 min]
% -- NARRATE
The base spends stopping clearance. Pressing a latch, the arm spends force
and deformation margin: small position errors against a stiff fixture
make large forces. Holding a load still, the machine does no mechanical
work yet must sustain stall current, so it spends heat and supply margin.
One task crosses all three.
% -- ENGAGE
Ask: ``The running machine opens the cage door. Which classes does that
one task cross?'' [Class 1 approaching, Class 2 at the latch, Class 3 if
it holds the door against the spring.]
% -- WARN
The classes name dominant constraints, not hardware categories.
% -- FLEX
[OPTIONAL] The humanoid example (a recovery step that needs peak torque,
current, and heat at once) previews the part-ending humanoid boxes.
}

\centering
\safeimg[width=0.95\textwidth,height=4.05cm,keepaspectratio]{three-classes.pdf}

\raggedright\footnotesize
The running machine crosses all three: stopping clearance in the aisle, contact force at the latch, winding heat and bus voltage while holding a load.

\bottomline{The machine classes are named for the margin each spends; one task can spend all three.}
\end{frame}

% =============================================================================
% SECTION 4: SCOPE TEST
% =============================================================================
\section{Scope Test}

\begin{frame}{Physical AI needs all three: learned, physical, delegated}
\note{
[3 min]
% -- NARRATE
Define each circle. Learned: behavior emerges from training data, so
correctness must be established empirically over an operating domain; a
linear plant model from system identification, with provable stability
margins, does not count. Physical feedback: actions change the material
world, and those changes govern what the sensors read next. Delegated
authority: outputs reach execution without synchronous human approval of
each control cycle.
% -- ENGAGE
Ask: ``Name a system for each pairwise overlap.'' [Advisory: learned
toolpaths a technician inspects before loading them onto a mill. Digital
autonomy: a trading engine. Classical control: a PID speed governor.]
% -- FLEX
[CORE]
}

\begin{columns}[T]
\begin{column}{0.46\textwidth}
\safeimg[width=\textwidth,height=4.9cm,keepaspectratio]{scope-venn.pdf}
\end{column}
\begin{column}{0.51\textwidth}
\small
\textbf{1. Learned statistical model:} behavior comes from training data, not an inspectable closed-form specification.

\vspace{0.1cm}
\textbf{2. Consequential physical feedback:} actions change the world, and the change governs what the sensors see next.

\vspace{0.1cm}
\textbf{3. Delegated authority:} outputs reach actuation without synchronous human approval for each control cycle.

\end{column}
\end{columns}

\bottomline{Excluding the overlaps leaves the systems where statistical uncertainty collides directly with physical irreversibility.}
\end{frame}

\begin{frame}{Delegated authority survives a veto, not per-command approval}
\note{
[3 min]
% -- NARRATE
Two edge cases decide most arguments. Learned perception feeding a
hand-written planner still has delegated authority, if an unvetted
detection can trigger a swerve or a brake; the Tempe vehicle was such a
machine. And a deterministic veto, such as an analog brake relay driven
by a certified safety lidar, does not remove authority, because learned
proposals still govern all nominal motion up to the threshold.
% -- ENGAGE
Ask: ``Does a surgeon at the console remove the assistive policy's
authority?'' [No: the haptic forces act within 150 to 250 ms reaction
time, without the surgeon acknowledging each correction.]
% -- WARN
Students think adding safety checks takes a system out of scope. Only
screening every command before transmission does.
% -- FLEX
[CORE]
}

\small
\begin{tabular}{@{}L{0.42\textwidth}L{0.12\textwidth}L{0.38\textwidth}@{}}
  \toprule
  \textbf{Design} & \textbf{In scope?} & \textbf{Why} \\
  \midrule
  Learned perception feeds a hand-written planner & yes & a learned detection can trigger braking or a swerve \\
  Learned proposals plus an independent deterministic veto & yes & proposals still drive all nominal motion \\
  Learned haptic guidance on a surgical console & yes & forces act inside human reaction time \\
  A person screens and authorizes every trajectory & no & synchronous approval removes the authority \\
  \bottomrule
\end{tabular}

\bottomline{Evaluate the three criteria independently; a learned network need not write to motor registers to hold physical authority.}
\end{frame}

\begin{frame}{Your turn: which machines are physical AI?}
\note{
[4 min] Exercise in pairs.
% -- NARRATE
Two minutes in pairs, then poll each line. Ask students to name the
missing criterion for every ``no.''
% -- WARN
Line 3 is the trap: the veto makes students say no. It stays in scope.
% -- FLEX
[CORE]
}

\small
For each machine, apply all three criteria. Which are physical AI, and which criterion is missing otherwise?
\begin{enumerate}\setlength{\itemsep}{2pt}
  \item A diffusion model plans CNC toolpaths; an engineer inspects and loads each one.
  \item A vision model highlights defects on an assembly-line display.
  \item The running machine: a chunk policy drives the base, under a lidar-triggered brake relay.
  \item A PID governor holds a turbine at speed.
\end{enumerate}

\pause
\begin{mlsyscard}{crimson}
\footnotesize
\textbf{1.} No: advisory, the engineer holds authority. \textbf{2.} No: assists a person, closes no physical loop. \textbf{3.} \textbf{Yes:} the veto does not remove delegated authority. \textbf{4.} No: classical control, no learned model.
\end{mlsyscard}
\end{frame}

% =============================================================================
% SECTION 5: FIVE LEVELS
% =============================================================================
\section{Five Levels}

\begin{frame}{The machine in five levels, from learning down to physics}
\note{
[3 min]
% -- NARRATE
Top to bottom. The Brain: learned models that propose goals, trajectories,
or torques, at 2 Hz for intent, 20 Hz chunk renewal, 40 Hz setpoints. The
Nervous System: buses, schedulers, and the permission path, a 1 kHz loop
in one implementation. The Body: drives, transmissions, windings, with
current loops at 10 to 25 kHz. The Boundary: register writes and sensor
captures, event-timed. The World: continuous. Governance wraps all five
on human timescales; it is not a sixth level.
% -- ENGAGE
Ask: ``The list is by abstraction, not by the path a command travels.
Trace the command path.'' [Nervous System, then the Boundary at a
register write, then current in the Body's windings.]
% -- WARN
The rates are this reference design's, not physical constants. What is
fundamental is the gap at the proposal boundary.
% -- FLEX
[CORE]
}

\begin{columns}[T]
\begin{column}{0.60\textwidth}
\safeimg[width=\textwidth,height=5.2cm,keepaspectratio]{five-levels.pdf}
\end{column}
\begin{column}{0.37\textwidth}
\footnotesize
\textbf{Typical placement:} proposers on an accelerator host, the permission path on an independent real-time core, current loops on the motor drives.

\vspace{0.15cm}
The \textbf{Boundary} gathers both edges of the machine: register writes that energize windings, and sensor captures.
\end{column}
\end{columns}

\bottomline{Distributing responsibility across levels lets a failure be diagnosed by the level that failed, not blamed on the network weights.}
\end{frame}

\begin{frame}{Definition: the proposal boundary}
\note{
[3 min]
% -- NARRATE
The definition has three load-bearing words: independent, bounded latency,
and alone. The permission path checks fresh evidence, the stopping
envelope (the states from which the machine can still halt before
collision), and a fallback shown feasible from the current state. On the
running machine these checks are an unrated layer above separate certified
protective functions; the Nervous System chapter develops that split.
% -- ENGAGE
Ask: ``A reinforcement-learning policy emits joint torques directly. Does
it hold actuator-write authority?'' [No: its output is still a proposal
until the permission path admits it.]
% -- WARN
The pitfall students will build: a clamp on the model's output. A clamp
does nothing when the model hangs or falls silent.
% -- FLEX
[CORE]
}

\begin{mlsyscard}{crimson}
\small The \textbf{proposal boundary} is the line between a learned model, which may only \emph{propose} an action, and the \textbf{permission path}: an independent mechanism with bounded latency, which the model cannot delay or corrupt, and which alone may let the action reach an actuator.
\end{mlsyscard}

\vspace{0.1cm}
\small
\begin{columns}[T]
\begin{column}{0.48\textwidth}
\textbf{Admission checks every proposal against} fresh sensory evidence, the stopping envelope, and a fallback shown feasible from the current state.
\end{column}
\begin{column}{0.48\textwidth}
\textbf{Pitfall:} treating the permission path as an output clamp. It must take control when the proposer hangs, misses deadlines, or falls silent.
\end{column}
\end{columns}

\bottomline{A learned model holds no authority by itself; its output stays a proposal until an independent path admits it.}
\end{frame}

\begin{frame}{The proposal boundary decides who; the causal boundary, where}
\note{
[2 min]
% -- NARRATE
Students will conflate the two boundaries all term. Put them side by
side. The proposal boundary is not a level; it is the line between the
Brain and the Nervous System. The causal boundary is the actuation half
of the Boundary level.
% -- ENGAGE
Ask: ``A spoofed packet on the CAN bus commands the drive directly. Which
boundary did it bypass?'' [The proposal boundary: it reached the causal
boundary without passing the permission path.]
% -- FLEX
[CORE]
}

\begin{columns}[T]
\begin{column}{0.48\textwidth}
\begin{mlsyscard}{bookblue}
\small\textbf{Proposal boundary}\\[2pt]
\footnotesize Between the Brain and the Nervous System. Decides \textbf{who may approve} an action. Above it, components may only propose.
\end{mlsyscard}
\end{column}
\begin{column}{0.48\textwidth}
\begin{mlsyscard}{bookbrown}
\small\textbf{Causal boundary}\\[2pt]
\footnotesize At the register write that energizes a winding. Marks \textbf{where} an admitted command alters the plant's state.
\end{mlsyscard}
\end{column}
\end{columns}

\vspace{0.2cm}
\small
The Nervous System bridges the two: it verifies proposals against hard constraints before it energizes the drives. In control theory this pattern is \textbf{runtime assurance}: the Simplex architecture pairs a complex controller with a certified baseline and decision logic that chooses which one drives the plant.

\bottomline{Keep the two boundaries distinct: one is a line of authority, the other is a physical interface.}
\end{frame}

\begin{frame}{Independence must extend to the permission path's sensing}
\note{
[2 min]
% -- NARRATE
Separate processors are not enough. If the monitor validates a trajectory
using the same learned perception that produced it, one perception error
corrupts both. The lease protects against a stalled proposer, not against
a timely, confidently wrong one.
% -- ENGAGE
Ask: ``What would a sensing path with a distinct failure mode look like on
the running machine?'' [A certified scanner, wheel encoders, an IMU:
anything that does not pass through the learned perception stack.]
% -- WARN
Students equate ``independent'' with ``different chip.'' Independence is
about failure modes, not silicon.
% -- FLEX
[OPTIONAL]
}

\small
\begin{columns}[T]
\begin{column}{0.48\textwidth}
\begin{mlsyscard}{bookred}
\footnotesize\textbf{Shared perception, shared failure}\\[2pt]
A monitor that checks a proposal with the same learned perception that produced it fails in the same way at the same moment. A perception error yields \emph{timely}, unsafe proposals that no lease catches.
\end{mlsyscard}
\end{column}
\begin{column}{0.48\textwidth}
\begin{mlsyscard}{bookgreen}
\footnotesize\textbf{Blend, do not only reject}\\[2pt]
A permission path can project a proposal onto its constraints or blend in corrections, keeping task progress. Evaluation must then separate the proposer alone, the supervised machine, and the hazards the fallback itself adds.
\end{mlsyscard}
\end{column}
\end{columns}

\bottomline{The sensing that grants execution authority must have failure modes distinct from the learned models it supervises.}
\end{frame}

% =============================================================================
% SECTION 6: TEMPE
% =============================================================================
\section{Tempe}

\begin{frame}{Tempe, 2018: the responsibilities failed together}
\note{
[3 min]
% -- NARRATE
On a March night in 2018, a modified Volvo XC90 running Uber ATG's
developmental automated driving system struck and killed a pedestrian
pushing a bicycle across a road in Tempe, Arizona. Show the sensor suite
from the NTSB report: roof lidar, radars, cameras, and compute racks in
the cargo bay. The system detected her 5.6 s before impact.
% -- WARN
Be careful with blame. The NTSB's probable cause was the operator's
distraction. The design choices we trace explain why one unmonitored
person had become the last line of defense. The record does not show
that any prescribed architecture would have prevented the crash.
% -- FLEX
[CORE] Treat a fatal case with care in the room.
}

\begin{columns}[T]
\begin{column}{0.55\textwidth}
\safeimg[width=\textwidth,height=4.0cm,keepaspectratio]{ntsb-tempe-vehicle.jpg}\\[-2pt]
{\scriptsize\textcolor{midgray}{Test vehicle sensor locations. Source: NTSB, public domain.}}
\end{column}
\begin{column}{0.42\textwidth}
\small
The automated driving system \textbf{first detected} the pedestrian \textbf{5.6 s} before impact.

\vspace{0.15cm}
Detection alone did not yield a correct path prediction. The failure ran through perception and path prediction, a response rule that withheld braking, and an operator the design relied on as its last countermeasure.
\end{column}
\end{columns}

\bottomline{The collision does not reduce to an inaccurate model; the machine around the model failed.}
\end{frame}

\begin{frame}{Each relabeling discarded the evidence of motion}
\note{
[3 min]
% -- NARRATE
The system kept tracking the object, but its classification changed among
vehicle, bicycle, and other. When it changed, the path predictor did not
use the object's earlier locations, and for an ``other'' object it could
assume a static path. Successive positions are the evidence that something
is crossing the lane.
% -- WARN
The NTSB record says only that earlier locations were omitted. Do not
describe a Kalman-filter reset or a zero-velocity start; the report does
not.
% -- ENGAGE
Ask: ``Which demand on the permission path does this violate?'' [Fresh
evidence: evidence already in hand was discarded.]
% -- FLEX
[CORE]
}

\begin{columns}[T]
\begin{column}{0.48\textwidth}
\small
\textbf{What the record shows}
\begin{itemize}\setlength{\itemsep}{2pt}
  \item The object stayed tracked.
  \item Its class changed among \emph{vehicle}, \emph{bicycle}, and \emph{other}.
  \item After each change, the predictor omitted earlier locations; for \emph{other} it could assume a static path.
\end{itemize}
\end{column}
\begin{column}{0.48\textwidth}
\begin{mlsyscard}{bookred}
\small\textbf{Consequence}\\[2pt]
\footnotesize Each relabeling restarted the crossing pedestrian as an object with no recorded history of crossing. Path predictions did not place her in the vehicle's lane until \textbf{1.2 s} before impact.
\end{mlsyscard}
\end{column}
\end{columns}

\bottomline{A changing semantic label must not silently discard the evidence of motion that successive positions carry.}
\end{frame}

\begin{frame}{One second of suppression spent 18.7 m with no brake}
\note{
[4 min]
% -- NARRATE
At 1.2 s before impact, at 43.2 mph, the system recognized an emergency
and entered a one-second action-suppression period, meant to reduce
responses to false alarms. At 0.2 s, at 40.5 mph, it planned a gradual
slowdown and sounded an alert; it did not request emergency braking. The
operator disengaged at 0.02 s. Compute live: the average of the two
recorded speeds is 41.85 mph, about 18.7 m/s, so about 18.7 m in 1 s.
% -- ENGAGE
Ask: ``Speed still fell during that second. Why?'' [A slowdown for a
planned right turn was already under way; the suppression withheld the
emergency response.]
% -- WARN
Do not compute a counterfactual stop. The report does not establish when
a different design would have begun braking.
% -- FLEX
[CORE]
}

\centering
\safeimg[width=0.97\textwidth,height=3.6cm,keepaspectratio]{tempe-timeline.pdf}

\raggedright\footnotesize
The first term of the stopping budget, read off an accident record: $\tfrac{1}{2}(43.2 + 40.5)$ mph $\times$ 1 s $\approx$ \textbf{18.7 m}. The system capped its braking at a chosen \textbf{7 m/s$^2$ (0.71 g)}; police found the vehicle could brake at about \textbf{0.92 g (9.0 m/s$^2$)} there.

\bottomline{Detection, prediction, recognition, and braking are distinct events; reporting the first alone hides the failures between them.}
\end{frame}

\begin{frame}{A disabled redundancy and an unmonitored operator}
\note{
[3 min]
% -- NARRATE
Read the NTSB findings. Probable cause: the operator was visually
distracted by her phone. Contributing: Uber ATG's inadequate safety risk
assessment, ineffective oversight of operators, and no mechanism for
automation complacency, all from an inadequate safety culture. Uber ATG
had also deactivated the Volvo's factory collision warning and automatic
emergency braking over possible radar interference. Volvo simulations
the NTSB cites suggest that system might have prevented or mitigated the
crash.
% -- ENGAGE
Ask: ``Cruise (2023) is the opposite pole. Why?'' [Tempe failed before
impact; Cruise moved after impact, about 6.1 m at up to 7.7 mph, without
fresh evidence about the space beneath the vehicle.]
% -- FLEX
[CORE]
}

\begin{columns}[T]
\begin{column}{0.49\textwidth}
\begin{mlsyscard}{bookred}
\footnotesize\textbf{Tempe (2018), before impact}\\[2pt]
The factory emergency braking had been \textbf{deactivated}, removing a layer of redundancy. The design made the operator the primary emergency countermeasure, yet nothing monitored her attention: \textbf{authority rested on evidence nobody collected}.
\end{mlsyscard}
\end{column}
\begin{column}{0.47\textwidth}
\begin{mlsyscard}{bookviolet}
\footnotesize\textbf{Cruise (2023), after impact}\\[2pt]
After braking hard and striking a pedestrian thrown into its path, the vehicle began a pullover and moved about \textbf{6.1 m at up to 7.7 mph} with her beneath it; its sensing did not detect her there.
\end{mlsyscard}
\end{column}
\end{columns}

\vspace{0.1cm}
\footnotesize
A fallback is itself motion, so it needs the same permission check as any proposal.

\bottomline{Safety needs timely hazard assessment and a mitigation path even when perception fails; accountability needs the record of what the machine knew.}
\end{frame}

\begin{frame}{Discussion: Tempe through the five levels}
\note{
[4 min] Discussion in pairs.
% -- NARRATE
Two minutes in pairs, then take one answer per question. Expected: (1)
Brain-level perception and prediction discarded evidence; the Nervous
System's response rule withheld braking; the factory brake, an
independent path, had been removed. (2) No: a person cannot be a
bounded-latency permission path without a mechanism showing they are
attending, and human reaction alone is 150 to 250 ms. (3) Fresh evidence,
deterministic authority with a feasible fallback, and recorded telemetry.
% -- WARN
Steer away from ``the model should have been better.'' The lecture's
claim is that the machine around the model failed.
% -- FLEX
[CORE]
}

\centering
\vspace{0.2cm}
{\large The record shows classification churn, a one-second suppression,\\ a deactivated factory brake, and an unmonitored operator.}\\[0.4cm]
\raggedright\small
\begin{enumerate}\setlength{\itemsep}{4pt}
  \item Assign each failure to a level of the five-level machine.
  \item Was the human operator a valid permission path? What would have had to be true?
  \item Which independent path, had it stayed enabled, sat outside the learned proposer?
\end{enumerate}
\centering
{\small\textcolor{midgray}{Two minutes with a neighbor.}}
\end{frame}

% =============================================================================
% SECTION 7: FOUR LAWS
% =============================================================================
\section{Four Laws}

\begin{frame}{Safety, security, and accountability converge in one place}
\note{
[3 min]
% -- NARRATE
In enterprise systems these live in separate tiers: security at network
boundaries, reliability in retries and circuit breakers, accountability in
append-only logs. On a robot crossing an aisle, a dirty camera, spoofed
packets on the CAN bus, and a hallucinated clear path all end in the same
kinetic hazard. Admission control cannot just confirm who sent a command;
it must decide whether executing it now violates the machine's limits.
% -- ENGAGE
Ask: ``A command carries a valid cryptographic signature. Is it
admissible?'' [Not necessarily: credentials establish the right to
request, not physical admissibility in the current state.]
% -- FLEX
[CORE]
}

\centering
\safeimg[width=0.95\textwidth,height=3.8cm,keepaspectratio]{convergence.pdf}

\raggedright\footnotesize
Physical authority is \textbf{state-dependent}. A joint command of 2 rad/s is routine in free space and unacceptable within 2 mm of a hard stop or a coworker.

\bottomline{Because momentum cannot be rolled back, every threat to a physical machine reduces to three demands on the permission path.}
\end{frame}

\begin{frame}{A slick floor invalidates every envelope built on 2 m/s$^2$}
\note{
[3 min]
% -- NARRATE
The card's 2 m/s$^2$ rests on a floor friction coefficient of at least
0.204. Spilled hydraulic fluid can drop tire-floor friction from about 0.8
to 0.2, an illustrative pair. When traction caps deceleration, braking
distance grows in inverse proportion to friction: a vehicle braking at the
traction limit sees it quadruple. The camera cannot see a transparent
slick, so the planner stays confident.
% -- ENGAGE
Ask: ``Which signals could reveal the slick in time?'' [Wheel-slip
telemetry, disagreement between commanded and IMU acceleration, motor
current transients.]
% -- WARN
Detection after the clearance is spent supports harm mitigation, not
collision avoidance.
% -- FLEX
[CORE]
}

\begin{columns}[T]
\begin{column}{0.50\textwidth}
\small
\textbf{The premise:} credible deceleration of 2 m/s$^2$ assumes friction $\mu \ge 0.204$.

\vspace{0.1cm}
\textbf{The change:} a hydraulic spill drops friction from about 0.8 to 0.2 (illustrative). At the traction limit, braking distance $\propto 1/\mu$: \textbf{4$\times$ longer}.

\vspace{0.1cm}
\textbf{The result:} 0.2 falls below the premise, so every envelope and speed profile computed from 2 m/s$^2$ a moment earlier is invalid.
\end{column}
\begin{column}{0.46\textwidth}
\begin{mlsyscard}{bookgreen}
\footnotesize\textbf{Independent evidence from the plant}\\[2pt]
wheel-slip telemetry\\
commanded vs.\ IMU acceleration\\
motor-current transients\\[3pt]
Restrict motion with a response feasible under the updated traction estimate.
\end{mlsyscard}
\end{column}
\end{columns}

\bottomline{Model confidence is not proof of physical feasibility; the permission path needs fresh evidence from the plant itself.}
\end{frame}

\begin{frame}{The three demands the permission path must discharge}
\note{
[2 min]
% -- NARRATE
State the three demands, then map the incidents onto them. Tempe's
predictor discarded evidence it already had, and the independent factory
brake had been disabled. Cruise moved without fresh evidence about the
space beneath it. Both investigations depended on the third demand.
% -- WARN
Synchronized means hardware clocks (IEEE 1588 PTP), because investigators
must separate a planning failure, an admission error, and an actuator
breakdown. Command logs alone show intent, not what the plant did.
% -- FLEX
[CORE]
}

\footnotesize
\begin{tabular}{@{}L{0.32\textwidth}L{0.60\textwidth}@{}}
  \toprule
  \textbf{Demand} & \textbf{What it requires, and where it failed} \\
  \midrule
  \textbf{Fresh, authenticated physical state} & fresh enough to bound uncertainty, robust to tampering. Tempe discarded tracking history; Cruise moved blind to the space beneath it. \\
  \textbf{Deterministic authority to intervene} & power to veto or clamp and inject a feasible fallback, whatever the model's confidence. Tempe's factory brake was off. \\
  \textbf{Synchronized causal telemetry} & an immutable, clock-synchronized record of what was proposed, evaluated, dispatched, and done. Both cases were rebuilt from it. \\
  \bottomrule
\end{tabular}

\bottomline{The permission path is where safety, security, and accountability obligations are discharged.}
\end{frame}

\begin{frame}{Four bedrock laws, one for each part of the book}
\note{
[4 min]
% -- NARRATE
These four laws are the spine of the volume. Read each law with its
consequence. First: checks must complete before the machine spends its
clearance. Second: single-step accuracy on a benchmark cannot guarantee
closed-loop behavior. Third: fresh state, an envelope check, and an
executable fallback, with a permission path the model cannot delay,
starve, or corrupt. Fourth: an operational design domain, its
assumptions, and degradation when conditions drift outside it.
% -- WARN
Use the exact names. ``Proposal is not permission'' and ``permission must
be earned'' are maxims; keep them verbatim.
% -- FLEX
[CORE]
}

\footnotesize
\begin{tabular}{@{}L{0.28\textwidth}L{0.46\textwidth}L{0.18\textwidth}@{}}
  \toprule
  \textbf{Law} & \textbf{What it demands} & \textbf{Part} \\
  \midrule
  \textbf{1. Irreversibility} & software cannot roll back physical consequences; checks finish before clearance is spent & I: Machine Anatomy \\
  \textbf{2. Endogenous data} & each action shapes the next observation; errors compound over the horizon & II: Teaching \\
  \textbf{3. Proposal is not permission} & an independent path admits every command, with a ready fallback & III: Running \\
  \textbf{4. Permission must be earned} & operate only where evidence supports performance under stated conditions & IV: Governing \\
  \bottomrule
\end{tabular}

\bottomline{Only the first is a physical law; the other three are design principles, elevated because violating them carries severe, predictable risk.}
\end{frame}

\begin{frame}{The laws distill established ideas and documented failures}
\note{
[3 min]
% -- NARRATE
The first law follows from mechanics and thermodynamics: work on a plant
transfers energy and raises entropy. The other three have established
names, and students should use those names when searching the literature.
Read the failure column: Tempe's 18.7 m; compounding error as a learning
result; Therac-25, where software took over protection that hardware
interlocks had given, and Mars Pathfinder's priority inversion.
% -- ENGAGE
Ask: ``Which row's failure is not an accident?'' [The second: it is a
theorem about imitation learning.]
% -- FLEX
[OPTIONAL] Fresh evidence (age of information, Kaul et al. 2012) is the
input to every permission check; mention it if time allows.
}

\scriptsize
\begin{tabular}{@{}L{0.19\textwidth}L{0.12\textwidth}L{0.27\textwidth}L{0.32\textwidth}@{}}
  \toprule
  \textbf{Law} & \textbf{Kind} & \textbf{Established name} & \textbf{Documented failure it distills} \\
  \midrule
  Irreversibility & physical & irreversible work; reaction-plus-braking stopping budget & Tempe (2018): 18.7 m with emergency braking withheld \\
  Endogenous data & design principle & covariate shift; compounding error in imitation learning (Ross et al.\ 2011) & a learning result: a cloned policy's errors compound with the horizon \\
  Proposal is not permission & design principle & runtime assurance, Simplex; safety filter (Sha 2001; Hsu et al.\ 2024) & Therac-25; Mars Pathfinder (1997) priority inversion \\
  Permission must be earned & design principle & assurance case; operational design domain (UL 4600) & Tempe's unmonitored operator; Cruise's unverified pullover \\
  \bottomrule
\end{tabular}

\bottomline{Where the book's terms add precision, such as tying freshness to the stopping budget, it keeps its own phrasing; the classical names index the literature.}
\end{frame}

\begin{frame}{The four laws meet at the proposal boundary}
\note{
[3 min]
% -- NARRATE
Read the four lines as one argument: irreversibility says when the check
must finish; endogenous data says why offline metrics cannot replace it;
proposal is not permission says where it lives; permission must be earned
says within which conditions. Then the scaling point: a bigger model does
not relax any of this, because onboard compute draws on fixed memory
bandwidth, power, and heat, and latency on the stopping path is travel.
% -- ENGAGE
Ask: ``The permission path rejects a command. What must happen next?''
[An admissible fallback must be dispatched immediately; halting a thread
leaves momentum untouched.]
% -- FLEX
[CORE]
}

\small
\begin{itemize}\setlength{\itemsep}{0pt}
  \item \textbf{Irreversibility:} checks complete before a command alters the plant.
  \item \textbf{Endogenous data:} offline metrics cannot certify the proposer; check it in the loop.
  \item \textbf{Proposal is not permission:} the check sits outside the model, on a path it cannot delay.
  \item \textbf{Permission must be earned:} admitted proposals stay within verified conditions.
\end{itemize}

\begin{mlsyscard}[top=0.15em,bottom=0.15em]{bookorange}
\footnotesize\textbf{Scale does not relax the laws.} Every added layer draws on fixed onboard bandwidth, power, and heat budgets; at the drive limit each millisecond on the stopping path costs 1.5 mm.
\end{mlsyscard}

\bottomline{A refusal means something only if a fallback the body can execute within its remaining clearance follows at once.}
\end{frame}

\begin{frame}{Each chapter adds a record; the chain meets at permission}
\note{
[3 min]
% -- NARRATE
A record is a machine-checkable artifact that later decisions can verify.
Today's stopping calculation is already a sketch of the first one, a
limit record, with six elements: nominal value, operating conditions,
uncertainty, margin, runtime detector, and provenance. Ours has only the
first two: the 762.9 mm stop, and payload, dry floor, stationary obstacle,
constant speed. Two strands run through the book: what the plant can do
and what the proposer has shown, and what the proposer proposes now.
They meet at the permission decision.
% -- WARN
A record documents a claim; it does not guarantee its physical truth.
% -- FLEX
[OPTIONAL] Name the records only as forward pointers; each is owned by
its chapter.
}

\footnotesize
\begin{tabular}{@{}L{0.20\textwidth}L{0.74\textwidth}@{}}
  \toprule
  \textbf{Part} & \textbf{Records it produces} \\
  \midrule
  I, Anatomy & limit record, handoff record \\
  II, Teaching & provenance record, policy manifest, evaluation record \\
  III, Running & observation contract, belief record, intent lease, trajectory record, enforcement record, placement record \\
  IV, Governing & authority record, fault record, release manifest, approval record \\
  \bottomrule
\end{tabular}

\vspace{0.15cm}
\begin{mlsyscard}{bookblue}
\footnotesize\textbf{Today's stop as a limit record.} Nominal value 762.9 mm; conditions: loaded, dry floor, stationary obstacle, constant speed. Still missing: quantified uncertainty, a verified margin, a runtime detector, and calibrated provenance.
\end{mlsyscard}

\bottomline{An ungrounded calculation becomes engineering evidence only when every field of its record is populated and backed by data.}
\end{frame}

% =============================================================================
% SECTION 8: WRAP-UP
% =============================================================================
\section{Wrap-Up}

\begin{frame}{Fallacies}
\note{
[3 min]
% -- NARRATE
Read each claim, ask for hands on ``true,'' then give the rebuttal.
% -- FLEX
[CORE]
}

\small
\textbf{Fallacy:} \textit{Safety constraints can be trained into the policy.}\\
Penalty terms suppress unsafe actions over the training distribution but guarantee nothing on out-of-distribution inputs. A loss optimizes expected performance; safety needs invariants that hold on every control cycle, so enforcement sits downstream of the policy on an independent permission path.

\vspace{0.25cm}
\textbf{Fallacy:} \textit{Physical actuation is reversible through later software compensation.}\\
A reverse motor command cannot unbend plastically deformed metal or recover the $I^2R$ heat dissipated in the windings. Detection after contact only truncates damage; checking trajectories against torque limits and kinetic-energy caps before dispatch is the primary defense.
\end{frame}

\begin{frame}{Pitfalls}
\note{
[3 min]
% -- NARRATE
Each pitfall treats time as if it were free. For the first: 605 ms at a
1 kHz loop is about 600 control cycles that no human checked.
% -- FLEX
[CORE]
}

\small
\textbf{Pitfall:} \textit{Treating human supervisory intervention as instantaneous.}\\
Simple human reaction takes 150 to 250 ms and grows under load; with console video and a handshake the running machine budgets an \textbf{in-loop takeover of 605 ms}. A 1 kHz loop executes hundreds of unchecked cycles in that time, so local force, velocity, and displacement envelopes must hold without a human.

\vspace{0.25cm}
\textbf{Pitfall:} \textit{Assuming software latency merely degrades performance gracefully.}\\
A manipulator that repeats its last forward-velocity command through a \textbf{100 ms} inference hiccup can hit an unmapped obstacle. Commands need bounded validity (the chunk lease), enforced by watchdogs, with a feasible fallback when it expires.
\end{frame}

\begin{frame}{What were the key ideas?}
\note{
[3 min] Retrieval practice.
% -- NARRATE
Timer on. Then ask three students for one item each and build the list on
the board. Compare with the next slide.
% -- FLEX
[CORE]
}

\centering
\vspace{0.8cm}
{\Large\bfseries Close your notes.}\\[0.5cm]
{\normalsize Write down \textbf{the elementary stopping budget}, \textbf{the three scope criteria},\\ \textbf{the five levels}, and \textbf{the four laws}.}\\[0.3cm]
{\normalsize 90 seconds, no peeking.}
\end{frame}

\begin{frame}{Key takeaways}
\note{
[2 min]
% -- NARRATE
One bullet per section. Return to the learning objectives and the
subtitle: cannot be recalled, and delay costs distance.
% -- FLEX
[CORE]
}

\footnotesize
\begin{itemize}\setlength{\itemsep}{0pt}
  \item \textbf{Beyond glass:} embodiment removes elastic time, rollback, and batching.
  \item \textbf{Causal boundary:} where a command first energizes a winding; past it, recovery is new motion.
  \item \textbf{Stopping budget:} $\dstop = v\,\taudelay + v^2/(2\abrake) \le \Dclear$: 200.4 + 562.5 = 762.9 mm at 1.5 m/s, and each millisecond is 1.5 mm; $v^* \approx 0.53$ m/s.
  \item \textbf{Scope test:} learned model, physical feedback, and delegated authority, all at once.
  \item \textbf{Five levels:} the proposal boundary sets who may approve; the causal boundary, where work begins.
  \item \textbf{Tempe:} discarded evidence, 18.7 m of withheld braking, a disabled brake, an unmonitored operator.
  \item \textbf{Four laws:} irreversibility (physical), endogenous data, proposal is not permission, permission must be earned.
\end{itemize}
\vfill\begin{mlsyscard}[top=0.15em,bottom=0.15em,fontupper=\footnotesize]{crimson}\textbf{Today's thesis.} Once a command becomes physical work it cannot be recalled, so delay costs \mbox{distance}.\end{mlsyscard}
\end{frame}

\begin{frame}{References}
\note{
[0 min] Reference list for further reading.
% -- FLEX
[OPTIONAL]
}

\mlsysrefsmall{Lee08}{Lee, ``Cyber Physical Systems: Design Challenges,'' ISORC 2008.}
\mlsysrefsmall{Sha01}{Sha, ``Using Simplicity to Control Complexity,'' IEEE Software 2001.}
\mlsysrefsmall{Ross11}{Ross, Gordon, and Bagnell, ``A Reduction of Imitation Learning \ldots\ to No-Regret Online Learning,'' AISTATS 2011.}
\mlsysrefsmall{Leveson93}{Leveson and Turner, ``An Investigation of the Therac-25 Accidents,'' Computer 1993.}
\mlsysrefsmall{Leveson11}{Leveson, \emph{Engineering a Safer World}, MIT Press 2011.}
\mlsysrefsmall{Hsu24}{Hsu, Hu, and Fisac, ``The Safety Filter: A Unified View of Safety-Critical Control,'' Annu.\ Rev.\ 2024.}
\mlsysrefsmall{NTSB19}{NTSB, \emph{Collision Between Vehicle Controlled by Developmental ADS and Pedestrian, Tempe}, HAR-19/03, 2019.}
\mlsysrefsmall{UL4600}{UL Standards \& Engagement, \emph{ANSI/UL 4600: Evaluation of Autonomous Products}, 2023.}
\mlsysrefsmall{Epoch26}{Epoch AI, \emph{Data on AI Models}, 2026.}
\end{frame}

\begin{frame}{Next lecture: which physical budget binds first?}
\note{
[1 min]
% -- LINK
Today we priced one margin, stopping clearance, with illustrative inputs.
Next time we profile the plant itself.
% -- FLEX
[CORE]
}

\centering
\vspace{0.4cm}
{\large A permission decision needs a physical limit against which to judge the proposed action.}\\[0.4cm]
{\normalsize Measurement freshness, stopping distance, deliverable torque, winding heat, and supply voltage: for a given maneuver, which binds first, under what stress, and on what evidence?}\\[0.6cm]
\begin{mlsyscard}{crimson}
\centering\normalsize \textbf{The Physical Body} turns each budget into a limit record, the six-element form today's stopping calculation only sketched.
\end{mlsyscard}
\end{frame}

\end{document}
